\documentclass[11pt,a4paper]{article}
\usepackage[utf8]{inputenc}
\usepackage[T1]{fontenc}
\usepackage{geometry}
\usepackage{xcolor}
\usepackage{hyperref}
\usepackage{titlesec}

\geometry{margin=1in}
\definecolor{tumblue}{RGB}{0,101,189}

\hypersetup{
    colorlinks=true,
    linkcolor=tumblue,
    urlcolor=tumblue,
    pdftitle={Cover Letter - Ismail AISSAOUI},
}

\begin{document}

\noindent \textbf{Ismail AISSAOUI} \\
State Engineer in Artificial Intelligence (ESI-SBA) \\
Email: ismail.aissaoui@example.com \\
Phone: [Your Phone Number] \\
LinkedIn: [Your Link]

\bigskip

\begin{flushright}
    Professor Wolfgang Polifke \\
    \textbf{Technical University of Munich (TUM)} \\
    Professorship of Thermofluid Dynamics \\
    Munich, Germany \\
    \medskip
    May 5, 2026
\end{flushright}

\bigskip

\noindent \textbf{Subject: Application for Doctoral Research - Physics-Informed Digital Twins}

\bigskip

Dear Professor Polifke,

I am writing to formally submit my application for a doctoral position under your supervision at the Technical University of Munich. As an AI Engineer specializing in hybrid physics-data modeling, I have closely followed your pioneering research on Physics-Informed Deep Learning for thermofluid dynamics, and I am eager to contribute to this field.

Currently, I am concluding my state engineering thesis at **Sonatrach**, where I have been tasked with building a **Generative Digital Twin** for a fleet of LM2500 gas turbines. My research address the "reality gap" by embedding thermodynamic constraints directly into neural network architectures using **PINN** loss functions. I have found that this hybrid approach is essential for maintaining physical consistency in predictive maintenance scenarios where data can be sparse or noisy.

I am particularly interested in your work on modeling turbulent flows and combustion instabilities. My experience with **Graph Neural Networks (GNNs)** for modeling structural dependencies, combined with my proficiency in sequential modeling using **Mamba-SSM**, allows me to approach thermofluid problems from a high-performance system perspective. I am keen to investigate how these state-of-the-art architectures can be integrated into your thermofluid models to improve extrapolation capabilities in extreme operating conditions.

TUM is my first choice for a PhD due to its world-class research in Engineering and Informatics. I am a highly motivated and autonomous researcher, and I believe my industrial background at Sonatrach has prepared me to tackle the rigorous challenges of your laboratory.

I have attached my academic transcripts, my CV (noting that ESI-SBA is an H+ recognized institution), and a research proposal for your review. I would be honored to discuss my potential contributions to your group during an interview.

Thank you for your time and consideration.

Sincerely,

\bigskip

\noindent Ismail AISSAOUI

\end{document}
